% ============================================================
% Section 144: SRH复合与本征硅过剩载流子寿命
% ============================================================
\section{半导体物理：SRH复合与本征硅过剩载流子寿命}
\label{sec:srh-lifetime}

\begin{modelbox}[题目：本征硅中复合中心位于禁带中央时的过剩载流子寿命]
把一种复合中心杂质掺入本征硅中，如果它的能级位置在禁带中央，试证明小注入时过剩载流子寿命 $\tau=\tau_n+\tau_p$。
\end{modelbox}

\begin{tipbox}[考点快速分析]
\begin{itemize}
    \item \textbf{SRH复合理论}：单一复合中心能级的净复合率公式
    \item \textbf{本征条件}：$n_0=p_0=n_i$，$E_t=E_i$ 时 $n_1=p_1=n_i$
    \item \textbf{小注入}：$\Delta n=\Delta p\ll n_i$
    \item \textbf{少子寿命定义}：$\tau_{p0}=1/(N_tr_p)$，$\tau_{n0}=1/(N_tr_n)$
\end{itemize}
\end{tipbox}

\begin{derivationbox}[标准解题过程]
\textbf{1. SRH复合理论}

单一复合中心能级 $E_t$ 引起的净复合率：
\[
R=\frac{N_tr_nr_p(np-n_i^2)}{r_n(n+n_1)+r_p(p+p_1)},
\]
其中 $n_1=n_ie^{(E_t-E_i)/k_BT}$，$p_1=n_ie^{(E_i-E_t)/k_BT}$。

\textbf{2. 本征硅且 $E_t=E_i$}

对于本征硅，$n_0=p_0=n_i$。复合中心位于禁带中央，$E_t=E_i$，故 $n_1=p_1=n_i$。

非平衡小注入下：$n=n_i+\Delta n$，$p=n_i+\Delta p$，$\Delta n=\Delta p$，$\Delta p\ll n_i$。

分子：
\[
np-n_i^2=(n_i+\Delta n)(n_i+\Delta p)-n_i^2\approx2n_i\Delta p.
\]

分母：
\[
r_n(n+n_1)+r_p(p+p_1)=r_n(2n_i+\Delta p)+r_p(2n_i+\Delta p).
\]

\textbf{3. 过剩载流子寿命}

净复合率：
\[
R\approx\frac{N_tr_nr_p\cdot2n_i\Delta p}{r_n(2n_i+\Delta p)+r_p(2n_i+\Delta p)}.
\]

过剩载流子寿命 $\tau=\Delta p/R$：
\[
\tau=\frac{r_n(2n_i+\Delta p)+r_p(2n_i+\Delta p)}{N_tr_nr_p\cdot2n_i}.
\]

在小注入极限 $\Delta p\ll n_i$ 下：
\[
\tau\approx\frac{2n_i(r_n+r_p)}{N_tr_nr_p\cdot2n_i}=\frac{1}{N_tr_p}+\frac{1}{N_tr_n}.
\]

定义 $\tau_{p0}=1/(N_tr_p)$（强n型中少子空穴寿命），$\tau_{n0}=1/(N_tr_n)$（强p型中少子电子寿命），则
\begin{equation}
\boxed{\tau=\tau_{p0}+\tau_{n0}.}
\end{equation}
\end{derivationbox}

\begin{formulabox}[答案汇总]
\begin{center}
\begin{tabular}{ll}
\toprule
\textbf{结论} & \textbf{结果} \\
\midrule
本征硅中过剩载流子寿命 & $\tau=\tau_{p0}+\tau_{n0}$ \\
物理意义 & 寿命是两种少子寿命之和，复合过程受限于俘获较慢的一方 \\
\bottomrule
\end{tabular}
\end{center}
\end{formulabox}

\begin{insightbox}[物理图像]
\begin{itemize}
    \item SRH理论是描述通过深能级复合中心产生-复合过程的经典模型，是分析半导体少数载流子寿命的基础。
    \item 本征材料中因 $n_0=p_0=n_i$，且复合中心位于禁带中央，导致电子和空穴俘获的对称贡献，寿命为两者之和。
    \item 在非本征（掺杂）材料中，多子浓度远大于少子，分母中多子项占主导，寿命退化为相应的少子寿命。
\end{itemize}
\end{insightbox}
